\IfFileExists{stacks-project.cls}{%
\documentclass{stacks-project}
}{%
\documentclass{amsart}
}

% For dealing with references we use the comment environment
\usepackage{verbatim}
\newenvironment{reference}{\comment}{\endcomment}
%\newenvironment{reference}{}{}
\newenvironment{slogan}{\comment}{\endcomment}
\newenvironment{history}{\comment}{\endcomment}

% For commutative diagrams we use Xy-pic
\usepackage[all]{xy}

% We use 2cell for 2-commutative diagrams.
\xyoption{2cell}
\UseAllTwocells

% We use multicol for the list of chapters between chapters
\usepackage{multicol}

% This is generally recommended for better output
\usepackage{lmodern}
\usepackage[T1]{fontenc}

% For cross-file-references
\usepackage{xr-hyper}

% Package for hypertext links:
\usepackage{hyperref}

% For any local file, say "hello.tex" you want to link to please
% use \externaldocument[hello-]{hello}
\externaldocument[introduction-]{introduction}
\externaldocument[conventions-]{conventions}
\externaldocument[sets-]{sets}
\externaldocument[categories-]{categories}
\externaldocument[topology-]{topology}
\externaldocument[sheaves-]{sheaves}
\externaldocument[sites-]{sites}
\externaldocument[stacks-]{stacks}
\externaldocument[fields-]{fields}
\externaldocument[algebra-]{algebra}
\externaldocument[brauer-]{brauer}
\externaldocument[homology-]{homology}
\externaldocument[derived-]{derived}
\externaldocument[simplicial-]{simplicial}
\externaldocument[more-algebra-]{more-algebra}
\externaldocument[smoothing-]{smoothing}
\externaldocument[modules-]{modules}
\externaldocument[sites-modules-]{sites-modules}
\externaldocument[injectives-]{injectives}
\externaldocument[cohomology-]{cohomology}
\externaldocument[sites-cohomology-]{sites-cohomology}
\externaldocument[dga-]{dga}
\externaldocument[dpa-]{dpa}
\externaldocument[sdga-]{sdga}
\externaldocument[hypercovering-]{hypercovering}
\externaldocument[schemes-]{schemes}
\externaldocument[constructions-]{constructions}
\externaldocument[properties-]{properties}
\externaldocument[morphisms-]{morphisms}
\externaldocument[coherent-]{coherent}
\externaldocument[divisors-]{divisors}
\externaldocument[limits-]{limits}
\externaldocument[varieties-]{varieties}
\externaldocument[topologies-]{topologies}
\externaldocument[descent-]{descent}
\externaldocument[perfect-]{perfect}
\externaldocument[more-morphisms-]{more-morphisms}
\externaldocument[flat-]{flat}
\externaldocument[groupoids-]{groupoids}
\externaldocument[more-groupoids-]{more-groupoids}
\externaldocument[etale-]{etale}
\externaldocument[chow-]{chow}
\externaldocument[intersection-]{intersection}
\externaldocument[pic-]{pic}
\externaldocument[weil-]{weil}
\externaldocument[adequate-]{adequate}
\externaldocument[dualizing-]{dualizing}
\externaldocument[duality-]{duality}
\externaldocument[discriminant-]{discriminant}
\externaldocument[derham-]{derham}
\externaldocument[local-cohomology-]{local-cohomology}
\externaldocument[algebraization-]{algebraization}
\externaldocument[curves-]{curves}
\externaldocument[resolve-]{resolve}
\externaldocument[models-]{models}
\externaldocument[functors-]{functors}
\externaldocument[equiv-]{equiv}
\externaldocument[pione-]{pione}
\externaldocument[etale-cohomology-]{etale-cohomology}
\externaldocument[proetale-]{proetale}
\externaldocument[relative-cycles-]{relative-cycles}
\externaldocument[more-etale-]{more-etale}
\externaldocument[trace-]{trace}
\externaldocument[crystalline-]{crystalline}
\externaldocument[spaces-]{spaces}
\externaldocument[spaces-properties-]{spaces-properties}
\externaldocument[spaces-morphisms-]{spaces-morphisms}
\externaldocument[decent-spaces-]{decent-spaces}
\externaldocument[spaces-cohomology-]{spaces-cohomology}
\externaldocument[spaces-limits-]{spaces-limits}
\externaldocument[spaces-divisors-]{spaces-divisors}
\externaldocument[spaces-over-fields-]{spaces-over-fields}
\externaldocument[spaces-topologies-]{spaces-topologies}
\externaldocument[spaces-descent-]{spaces-descent}
\externaldocument[spaces-perfect-]{spaces-perfect}
\externaldocument[spaces-more-morphisms-]{spaces-more-morphisms}
\externaldocument[spaces-flat-]{spaces-flat}
\externaldocument[spaces-groupoids-]{spaces-groupoids}
\externaldocument[spaces-more-groupoids-]{spaces-more-groupoids}
\externaldocument[bootstrap-]{bootstrap}
\externaldocument[spaces-pushouts-]{spaces-pushouts}
\externaldocument[spaces-chow-]{spaces-chow}
\externaldocument[groupoids-quotients-]{groupoids-quotients}
\externaldocument[spaces-more-cohomology-]{spaces-more-cohomology}
\externaldocument[spaces-simplicial-]{spaces-simplicial}
\externaldocument[spaces-duality-]{spaces-duality}
\externaldocument[formal-spaces-]{formal-spaces}
\externaldocument[restricted-]{restricted}
\externaldocument[spaces-resolve-]{spaces-resolve}
\externaldocument[formal-defos-]{formal-defos}
\externaldocument[defos-]{defos}
\externaldocument[cotangent-]{cotangent}
\externaldocument[examples-defos-]{examples-defos}
\externaldocument[algebraic-]{algebraic}
\externaldocument[examples-stacks-]{examples-stacks}
\externaldocument[stacks-sheaves-]{stacks-sheaves}
\externaldocument[criteria-]{criteria}
\externaldocument[artin-]{artin}
\externaldocument[quot-]{quot}
\externaldocument[stacks-properties-]{stacks-properties}
\externaldocument[stacks-morphisms-]{stacks-morphisms}
\externaldocument[stacks-limits-]{stacks-limits}
\externaldocument[stacks-cohomology-]{stacks-cohomology}
\externaldocument[stacks-perfect-]{stacks-perfect}
\externaldocument[stacks-introduction-]{stacks-introduction}
\externaldocument[stacks-more-morphisms-]{stacks-more-morphisms}
\externaldocument[stacks-geometry-]{stacks-geometry}
\externaldocument[moduli-]{moduli}
\externaldocument[moduli-curves-]{moduli-curves}
\externaldocument[examples-]{examples}
\externaldocument[exercises-]{exercises}
\externaldocument[guide-]{guide}
\externaldocument[desirables-]{desirables}
\externaldocument[coding-]{coding}
\externaldocument[obsolete-]{obsolete}
\externaldocument[fdl-]{fdl}
\externaldocument[index-]{index}

% Theorem environments.
%
\theoremstyle{plain}
\newtheorem{theorem}[subsection]{Theorem}
\newtheorem{proposition}[subsection]{Proposition}
\newtheorem{lemma}[subsection]{Lemma}

\theoremstyle{definition}
\newtheorem{definition}[subsection]{Definition}
\newtheorem{example}[subsection]{Example}
\newtheorem{exercise}[subsection]{Exercise}
\newtheorem{situation}[subsection]{Situation}

\theoremstyle{remark}
\newtheorem{remark}[subsection]{Remark}
\newtheorem{remarks}[subsection]{Remarks}

\numberwithin{equation}{subsection}

% Macros
%
\def\lim{\mathop{\mathrm{lim}}\nolimits}
\def\colim{\mathop{\mathrm{colim}}\nolimits}
\def\Spec{\mathop{\mathrm{Spec}}}
\def\Hom{\mathop{\mathrm{Hom}}\nolimits}
\def\Ext{\mathop{\mathrm{Ext}}\nolimits}
\def\SheafHom{\mathop{\mathcal{H}\!\mathit{om}}\nolimits}
\def\SheafExt{\mathop{\mathcal{E}\!\mathit{xt}}\nolimits}
\def\Sch{\mathit{Sch}}
\def\Mor{\mathop{\mathrm{Mor}}\nolimits}
\def\Ob{\mathop{\mathrm{Ob}}\nolimits}
\def\Sh{\mathop{\mathit{Sh}}\nolimits}
\def\NL{\mathop{N\!L}\nolimits}
\def\CH{\mathop{\mathrm{CH}}\nolimits}
\def\proetale{{pro\text{-}\acute{e}tale}}
\def\etale{{\acute{e}tale}}
\def\QCoh{\mathit{QCoh}}
\def\Ker{\mathop{\mathrm{Ker}}}
\def\Im{\mathop{\mathrm{Im}}}
\def\Coker{\mathop{\mathrm{Coker}}}
\def\Coim{\mathop{\mathrm{Coim}}}

% Boxtimes
%
\DeclareMathSymbol{\boxtimes}{\mathbin}{AMSa}{"02}

%
% Macros for moduli stacks/spaces
%
\def\QCohstack{\mathcal{QC}\!\mathit{oh}}
\def\Cohstack{\mathcal{C}\!\mathit{oh}}
\def\Spacesstack{\mathcal{S}\!\mathit{paces}}
\def\Quotfunctor{\mathrm{Quot}}
\def\Hilbfunctor{\mathrm{Hilb}}
\def\Curvesstack{\mathcal{C}\!\mathit{urves}}
\def\Polarizedstack{\mathcal{P}\!\mathit{olarized}}
\def\Complexesstack{\mathcal{C}\!\mathit{omplexes}}
% \Pic is the operator that assigns to X its picard group, usage \Pic(X)
% \Picardstack_{X/B} denotes the Picard stack of X over B
% \Picardfunctor_{X/B} denotes the Picard functor of X over B
\def\Pic{\mathop{\mathrm{Pic}}\nolimits}
\def\Picardstack{\mathcal{P}\!\mathit{ic}}
\def\Picardfunctor{\mathrm{Pic}}
\def\Deformationcategory{\mathcal{D}\!\mathit{ef}}
\usepackage{cleveref}
\crefname{originallemma}{original lemma}{original lemmas}
\Crefname{originallemma}{Original lemma}{Original lemmas}
\makeatletter
\newlabel{varieties-lemma-geometrically-reduced}{{035X}{source}{}{sourcepointer-1}{}}
\newlabel{varieties-lemma-geometrically-reduced@cref}{{[originallemma][0][]035X}{[1][0][]source}{}{}{}}
\makeatother
\crefname{originallemma}{original lemma}{original lemmas}
\Crefname{originallemma}{Original lemma}{Original lemmas}
\makeatletter
\newlabel{algebra-lemma-pushdown-module}{{02M0}{source}{}{sourcepointer-2}{}}
\newlabel{algebra-lemma-pushdown-module@cref}{{[originallemma][0][]02M0}{[1][0][]source}{}{}{}}
\makeatother
\crefname{originallemma}{original lemma}{original lemmas}
\Crefname{originallemma}{Original lemma}{Original lemmas}
\makeatletter
\newlabel{algebra-lemma-Noetherian-field-extension}{{045I}{source}{}{sourcepointer-3}{}}
\newlabel{algebra-lemma-Noetherian-field-extension@cref}{{[originallemma][0][]045I}{[1][0][]source}{}{}{}}
\makeatother
\crefname{originallemma}{original lemma}{original lemmas}
\Crefname{originallemma}{Original lemma}{Original lemmas}
\makeatletter
\newlabel{algebra-lemma-make-separable}{{030R}{source}{}{sourcepointer-4}{}}
\newlabel{algebra-lemma-make-separable@cref}{{[originallemma][0][]030R}{[1][0][]source}{}{}{}}
\makeatother
\crefname{originallemma}{original lemma}{original lemmas}
\Crefname{originallemma}{Original lemma}{Original lemmas}
\makeatletter
\newlabel{algebra-lemma-localization-smooth-separable}{{037X}{source}{}{sourcepointer-5}{}}
\newlabel{algebra-lemma-localization-smooth-separable@cref}{{[originallemma][0][]037X}{[1][0][]source}{}{}{}}
\makeatother
\crefname{originallemma}{original lemma}{original lemmas}
\Crefname{originallemma}{Original lemma}{Original lemmas}
\makeatletter
\newlabel{algebra-lemma-pullback-module}{{02M1}{source}{}{sourcepointer-6}{}}
\newlabel{algebra-lemma-pullback-module@cref}{{[originallemma][0][]02M1}{[1][0][]source}{}{}{}}
\makeatother
\title{Multiplicities of components in a field extension base change are one or powers of the characteristic}
\author{The Stacks Project Authors}
\date{stacks-project:relative-cycles.tex:local-snapshot-20261012; GNU FDL 1.2 or later}
\begin{document}
\maketitle
\setcounter{subsection}{1}
\noindent\textit{Editorial note.} All delivered mathematical text is the author's own, copied from the pinned Stacks Project chapter file relative-cycles.tex; no statement or proof step was rewritten, reordered or removed. The counted span is the author environment and its own proof environment. References inside the span are handed to the assembler as navigation stubs that carry the canonical tag of the cited tag only; whether each of them is load-bearing is recorded on the stub itself.\par\medskip
\noindent\textit{Allowed standard prerequisites (editorial).}\begin{itemize}
\item Geometric reducedness of varieties over perfect fields under base change (Varieties, Lemma varieties-lemma-geometrically-reduced, canonical tag 035X). cited at proof lines 312-314. In characteristic zero k is perfect so the base change is reduced and every multiplicity is one.
\item Pushdown formula for lengths along a finite free ring map (Algebra, Lemma algebra-lemma-pushdown-module, canonical tag 02M0). cited at proof lines 337-339. Relates the length at a minimal prime of the smaller tensor product to the length at the unique minimal prime above it, scaled by the residue field degree and the rank.
\item Noetherianity of the tensor product with a field extension of finitely generated type (Algebra, Lemma algebra-lemma-Noetherian-field-extension, canonical tag 045I). cited at proof lines 318-319. Makes K tensor\_k L Noetherian, so its localizations at minimal primes are artinian local rings and lengths are defined.
\item Separable-algebraic splitting of field extensions (Algebra, Lemma algebra-lemma-make-separable, canonical tag 030R). cited at proof lines 345-347. Produces the subfield L/k'/k with L/k' separable and k'/k finite purely inseparable.
\item Separable extensions are localizations of smooth algebras (Algebra, Lemma algebra-lemma-localization-smooth-separable, canonical tag 037X). cited at proof lines 350-352. Makes the separable part a localization of a smooth k'-algebra.
\item Pullback formula for lengths along flat local homomorphisms with matched maximal ideals (Algebra, Lemma algebra-lemma-pullback-module, canonical tag 02M1). cited at proof lines 357-359. Transports the p-power length from the purely inseparable tensor product to the localized ring S.
\end{itemize}\medskip
\ifdefined\hypertarget\hypertarget{sourcepointer-1}{}\fi
\noindent\textit{Source reference.} Original lemma 035X: sources/stacks-chow/tags.txt, source lines 4144--4144. Cited at proof lines 312-314: multiplicity one in characteristic zero.\par\medskip
\ifdefined\hypertarget\hypertarget{sourcepointer-2}{}\fi
\noindent\textit{Source reference.} Original lemma 02M0: sources/stacks-chow/tags.txt, source lines 3427--3427. Cited at proof lines 337-339: descends the length statement along purely inseparable extensions.\par\medskip
\ifdefined\hypertarget\hypertarget{sourcepointer-3}{}\fi
\noindent\textit{Source reference.} Original lemma 045I: sources/stacks-chow/tags.txt, source lines 5437--5437. Cited at proof lines 318-319: makes K tensor\_k L Noetherian so the localizations at minimal primes are artinian.\par\medskip
\ifdefined\hypertarget\hypertarget{sourcepointer-4}{}\fi
\noindent\textit{Source reference.} Original lemma 030R: sources/stacks-chow/tags.txt, source lines 3958--3958. Cited at proof lines 345-346: produces the subfield L/k'/k.\par\medskip
\ifdefined\hypertarget\hypertarget{sourcepointer-5}{}\fi
\noindent\textit{Source reference.} Original lemma 037X: sources/stacks-chow/tags.txt, source lines 4216--4216. Cited at proof lines 350-352: makes L/k' a smooth localization.\par\medskip
\ifdefined\hypertarget\hypertarget{sourcepointer-6}{}\fi
\noindent\textit{Source reference.} Original lemma 02M1: sources/stacks-chow/tags.txt, source lines 3428--3428. Cited at proof lines 357-359: transports the length along R to S, finishing the proof.\par\medskip
\begin{lemma}
\label{lemma-multiplicities-field-extension}
Let $K/k$ be a field extension. Let $Z$ be an integral locally algebraic
scheme over $k$. The multiplicity $m_{Z', Z_K}$ of an irreducible
component $Z' \subset Z_K$ is $1$ or a power of the characteristic of $k$.
\end{lemma}

\begin{proof}
If the characteristic of $k$ is zero, then $k$ is perfect and
the multiplicity is always $1$ since $X_K$ is reduced by
Varieties, Lemma \ref{varieties-lemma-geometrically-reduced}.
Assume the characteristic of $k$ is $p > 0$.
Let $L$ be the function field of $Z$. Since $Z$ is locally algebraic
over $k$, the field extension $L/k$ is finitely generated.
The ring $K \otimes_k L$ is Noetherian
(Algebra, Lemma \ref{algebra-lemma-Noetherian-field-extension}).
Translated into algebra, we have to show that the length of the
artinian local ring $(K \otimes_k L)_\mathfrak q$
is a power of $p$ for every minimal prime ideal $\mathfrak q$.

\medskip\noindent
Let $L'/L$ be a finite purely inseparable extension, say of degree
$p^n$. Then $K \otimes_k L \subset K \otimes_k L'$ is a finite
free ring map of degree $p^n$ which induces a homeomorphism on
spectra and purely inseparable residue field extensions.
Hence for every minimal prime $\mathfrak q$ as above
there is a unique minimal prime
$\mathfrak q' \subset K \otimes_k L'$ lying over it and
$$
p^n \text{length}((K \otimes_k L)_\mathfrak q) =
[\kappa(\mathfrak q') : \kappa(\mathfrak q)]
\text{length}((K \otimes_k L')_{\mathfrak q'})
$$
by Algebra, Lemma \ref{algebra-lemma-pushdown-module} applied
to $M = (K \otimes_k L')_{\mathfrak q'} \cong
(K \otimes_k L)_{\mathfrak q}^{\oplus p^n}$.
Since $[\kappa(\mathfrak q') : \kappa(\mathfrak q)]$ is a power
of $p$ we conclude that it suffices to prove the
statement for $L'$ and $\mathfrak q'$.

\medskip\noindent
By the previous paragraph and Algebra, Lemma \ref{algebra-lemma-make-separable}
we may assume that we have a subfield $L/k'/k$ such that $L/k'$ is separable
and $k'/k$ is finite purely inseparable. Then $K \otimes_k k'$ is an
Artinian local ring. The argument of the preceding paragraph
(applied to $L = k$ and $L' = k'$) shows that $\text{length}(K \otimes_k k')$
is a power of $p$. Since $L/k'$ is the localization of a
smooth $k'$-algebra
(Algebra, Lemma \ref{algebra-lemma-localization-smooth-separable}).
Hence $S = (K \otimes_k L)_\mathfrak q$ is the localization of a smooth
$R = K \otimes_k k'$-algebra at a minimal prime.
Thus $R \to S$ is a flat local homomorphism of Artinian
local rings and $\mathfrak m_R S = \mathfrak m_S$. It
follows from Algebra, Lemma \ref{algebra-lemma-pullback-module} that
$\text{length}(K \otimes_k k') = \text{length}(R) =
\text{length}(S) = \text{length}((K \otimes_k L)_\mathfrak q)$
and the proof is finished.
\end{proof}

\end{document}
